\begin{definition}[The native plaquette walk and its physical sites]%
    \label{def:peps_torus_plaquette_flux_geometry}
    \lean{TNLean.PEPS.torusPlaquetteWalk,
        TNLean.PEPS.torusPlaquetteRegion,
        TNLean.PEPS.torusPlaquetteRegion_card}
    \leanok
    On a native torus whose two periods are at least three, let
    $v=(x,y)$. Define the plaquette walk $p_v$ and its physical region
    $R_v$ by
    \begin{align}
        p_v & :\ (x,y)\longrightarrow(x+1,y)
        \longrightarrow(x+1,y+1)\longrightarrow(x,y+1)
        \longrightarrow(x,y),                                     \\
        R_v & =\{(x,y),(x+1,y),(x+1,y+1),(x,y+1)\}, & |R_v| & =4.
        \label{eq:peps_plaquette_flux_walk}
    \end{align}
    Coordinates are periodic, so this definition includes plaquettes
    crossing either seam. It is the four-site region used in
    \cite[proof of Theorem~6.15]{Schuch2010PEPS}.
\end{definition}

\begin{definition}[A one-bond regular flux string]%
    \label{def:peps_torus_plaquette_flux_insertion}
    \lean{TNLean.PEPS.torusPlaquetteFluxInsertion,
        TNLean.PEPS.regularWalkHolonomy_torusPlaquetteFluxInsertion}
    \leanok
    \uses{def:peps_torus_plaquette_flux_geometry,
        thm:peps_regular_walk_holonomy}
    For $g\in G$, put directed transport $g$ on the first side of $p_v$
    and identity transport on every other bond. In the ordered-edge
    convention, the distinguished bond therefore carries $g$ when
    its orientation agrees with the walk and $g^{-1}$ otherwise.
    Denote this actual assignment by $u^{v,g}$. Its plaquette holonomy is
    \begin{align}
        \operatorname{Hol}_{u^{v,g}}(p_v) & =g.
        \label{eq:peps_plaquette_flux_holonomy}
    \end{align}
    This is the one-bond open regular string of
    \cite[Section~6.5 and proof of Theorem~6.15]{Schuch2010PEPS}.
\end{definition}
\begin{proof}\leanok
    The three other sides of the plaquette have at least one endpoint
    with vertical coordinate $y+1$, whereas both endpoints of the
    distinguished side have coordinate $y$. Hence their ordered edges
    are different. Their transports are the identity, and the
    orientation convention gives transport $g$ on the first step.
    The reverse-ordered product along the four steps is therefore $g$.
\end{proof}

\begin{theorem}[Four-site measurement of native plaquette holonomy]%
    \label{thm:peps_torus_plaquette_holonomy_measurement}
    \lean{TNLean.PEPS.IsGIsometric.exists_torusPlaquette_physicalMeasurement,
        TNLean.PEPS.IsGIsometric.exists_torusPlaquette_holonomy_cutMeasurement}
    \leanok
    \uses{def:peps_torus_plaquette_flux_geometry,
        thm:peps_closed_walk_class_physical_measurement,
        thm:peps_regular_boundary_route,
        thm:peps_regular_torus_incident_isometry,
        thm:peps_regular_open_region_projector}
    Let the two torus periods be at least three and let the original
    local tensor be regular $G$-isometric for a finite group $G$.
    At every vertex $v$, there is a single complete projective
    measurement $(Q_C)_C$, with one additional outcome $Q_\perp$, on
    the four original physical spins in $R_v$.
    Every operator is Hermitian and positive semidefinite, the operators
    satisfy $Q_iQ_j=\delta_{ij}Q_i$, and their sum is the identity.
    For every actual inserted bond assignment $u$ and every crossing
    configuration $\theta$, its twisted open-block column satisfies
    \begin{align}
        Q_C\psi_{u,\theta}
                               & =\mathbf1_{\{[\operatorname{Hol}_u(p_v)]=C\}}\psi_{u,\theta}, &
        Q_\perp\psi_{u,\theta} & =0.
        \label{eq:peps_plaquette_holonomy_measurement}
    \end{align}
    The same operators act on the actual globally contracted coefficient
    matrix $M_u$ across $R_v$ by
    \begin{align}
        Q_CM_u      & =\mathbf1_{\{[\operatorname{Hol}_u(p_v)]=C\}}M_u, &
        Q_\perp M_u & =0.
        \label{eq:peps_plaquette_arbitrary_closed_measurement}
    \end{align}
    This is the native regular torus version of the local conjugacy
    calculation in \cite[proof of Theorem~6.15]{Schuch2010PEPS}; see
    \cite{gap:rmp_peps_quantum_double_g_isometry} for its scope.
\end{theorem}
\begin{proof}\leanok
    The induced graph on the support of an actual walk is connected.
    Apply Theorem~\ref{thm:peps_closed_walk_class_physical_measurement}
    to $p_v$ in this induced graph. Inclusion of the induced walk in
    the ambient torus recovers $p_v$ and preserves its holonomy.
    Incident-edge coordinates preserve the original regular
    $G$-isometry, so the resulting operators act on the original four
    physical spins. All spanning-tree choices are made within this
    construction. Factoring the global coefficient matrix into its open
    block and complementary open block transfers the same eigenvalue
    from every open-block column to the contracted matrix.
\end{proof}

\begin{theorem}[Detection of a native one-bond flux string]%
    \label{thm:peps_torus_plaquette_flux_class_measurement}
    \lean{TNLean.PEPS.IsGIsometric.exists_torusPlaquette_fluxClass_measurement,
        TNLean.PEPS.IsGIsometric.exists_torusPlaquette_fluxClass_cutMeasurement}
    \leanok
    \uses{thm:peps_torus_plaquette_holonomy_measurement,
        def:peps_torus_plaquette_flux_insertion,
        thm:peps_regular_open_region_projector}
    Under the hypotheses of
    Theorem~\ref{thm:peps_torus_plaquette_holonomy_measurement}, one
    fixed complete projective measurement on the four spins of $R_v$
    detects the class $[g]$ of every one-bond flux string $u^{v,g}$.
    Its extra outcome annihilates every such actual open-block column.
    If $M_{v,g}$ is the physical coefficient matrix of the globally
    contracted inserted torus PEPS across the cut $R_v$, then
    \begin{align}
        Q_CM_{v,g}      & =\mathbf1_{\{[g]=C\}}M_{v,g}, &
        Q_\perp M_{v,g} & =0.
        \label{eq:peps_plaquette_flux_closed_measurement}
    \end{align}
    Thus the measurement detects $[g]$ with certainty on every nonzero
    such contracted state. This is the explicit one-bond string
    specialization of \cite[Theorem~6.15]{Schuch2010PEPS} on native
    regular tori of periods at least three.
\end{theorem}
\begin{proof}\leanok
    Substitute (\ref{eq:peps_plaquette_flux_holonomy}) into
    (\ref{eq:peps_plaquette_holonomy_measurement}). Factor the actual
    global coefficient matrix across the cut into its open block and
    complementary open block. Each column of the former has the same
    measured class $[g]$, so left multiplication by $Q_C$ multiplies
    the contracted matrix by $\mathbf1_{\{[g]=C\}}$.
\end{proof}

\begin{definition}[Flux strings attached to each plaquette side]%
    \label{def:peps_torus_plaquette_all_sides_insertion}
    \lean{TNLean.PEPS.torusPlaquetteSide,
        TNLean.PEPS.torusPlaquetteSideInsertion,
        TNLean.PEPS.regularWalkHolonomy_torusPlaquetteSideInsertion}
    \leanok
    \uses{def:peps_torus_plaquette_flux_geometry,
        thm:peps_regular_walk_holonomy}
    Enumerate the oriented sides of $p_v$ by $i=0,1,2,3$ in traversal
    order. For $g\in G$, let $u^{v,i,g}$ put directed transport $g$
    on side $i$ and identity transport on all other bonds. The ordered
    bond receives $g$ or $g^{-1}$ according to its orientation. Then
    \begin{align}
        \operatorname{Hol}_{u^{v,i,g}}(p_v) & =g.
        \label{eq:peps_plaquette_all_sides_holonomy}
    \end{align}
    This realizes the four attachment directions in the final paragraph
    of \cite[proof of Theorem~6.15]{Schuch2010PEPS}.
\end{definition}
\begin{proof}\leanok
    Opposite horizontal sides have distinct vertical coordinates, and
    opposite vertical sides have distinct horizontal coordinates.
    Adjacent sides also have distinct ordered endpoint pairs. Hence
    all four bonds are different. Exactly one factor in the actual
    four-step holonomy is $g$, and the other three are the identity.
\end{proof}

\begin{theorem}[A single flux measurement for all four attachment directions]%
    \label{thm:peps_torus_plaquette_all_sides_measurement}
    \lean{TNLean.PEPS.IsGIsometric.exists_torusPlaquette_allSides_cutMeasurement}
    \leanok
    \uses{thm:peps_torus_plaquette_holonomy_measurement,
        def:peps_torus_plaquette_all_sides_insertion}
    On native regular tori with both periods at least three, one fixed
    complete positive projective measurement on the four spins in $R_v$
    detects the class of every explicitly constructed string
    $u^{v,i,g}$, for all four attachment sides and all $g\in G$.
    Writing $M_{v,i,g}$ for its actual globally contracted coefficient
    matrix, one has
    \begin{align}
        Q_CM_{v,i,g}      & =\mathbf1_{\{[g]=C\}}M_{v,i,g}, &
        Q_\perp M_{v,i,g} & =0.
        \label{eq:peps_plaquette_all_sides_measurement}
    \end{align}
    The operators are independent of both $i$ and $g$. This is the
    explicit four-direction specialization of the final observation in
    \cite[proof of Theorem~6.15]{Schuch2010PEPS}; it does not assert
    identification with a parent-Hamiltonian excitation space.
\end{theorem}
\begin{proof}\leanok
    Substitute (\ref{eq:peps_plaquette_all_sides_holonomy}) into
    (\ref{eq:peps_plaquette_arbitrary_closed_measurement}).
\end{proof}

\begin{corollary}[Orthogonality of distinct measured plaquette classes]%
    \label{cor:peps_torus_plaquette_class_orthogonality}
    \lean{TNLean.PEPS.IsGIsometric.torusPlaquette_cutMatrix_orthogonal_of_class_ne}
    \leanok
    \uses{thm:peps_torus_plaquette_holonomy_measurement}
    Under the regular native torus hypotheses of
    Theorem~\ref{thm:peps_torus_plaquette_holonomy_measurement}, let
    $u,w$ be two actual bond assignments. If their holonomies around
    $p_v$ have distinct conjugacy classes, then their globally contracted
    coefficient matrices satisfy
    \begin{align}
        M_u^\dagger M_w & =0.
        \label{eq:peps_plaquette_class_orthogonality}
    \end{align}
    Thus every physical coefficient column of one is orthogonal to every
    coefficient column of the other, without a normalization hypothesis.
    This is a consequence of the local class measurement in
    \cite[Theorem~6.15]{Schuch2010PEPS}.
\end{corollary}
\begin{proof}\leanok
    Let $Q$ be the measurement projector for the class of
    $\operatorname{Hol}_u(p_v)$. Then $QM_u=M_u$ and $QM_w=0$.
    Since $Q$ is Hermitian,
    $M_u^\dagger M_w=M_u^\dagger Q M_w=0$.
\end{proof}
